\documentclass[12pt]{article}
\usepackage[utf8]{inputenc}
\usepackage[T1]{fontenc}
\usepackage{amsmath, amssymb, geometry}
\usepackage{enumitem}
\geometry{margin=1in}
\setlength{\parindent}{0pt}
\renewcommand{\arraystretch}{1.2}

\title{Solutions -- Quiz 5\\ \vspace{5pt} \Large ECON 308 -- International Economic Relations}
\author{Instructor: Muhebullah Karimzada}
\date{December 02,  2025}
\usepackage{comment}

\begin{document}
\maketitle

\section*{Multiple-Choice Questions – Answer Key and Explanations}

\begin{enumerate}[label=\textbf{\arabic*.}]

\item \textbf{Correct answer: B} \\
A \textit{specific tariff} is a fixed amount of money per physical unit of the imported good (e.g.\ \$10 per ton).  
Option A describes an \emph{ad valorem} tariff (percentage of value).

\vspace{0.5em}

\item \textbf{Correct answer: C} \\
For a \textit{small country}, its trade policy does not affect world prices, so there is \emph{no terms-of-trade gain}.  
A tariff only creates efficiency losses (production and consumption distortions), so the net welfare effect is purely a \textbf{deadweight loss}.

\vspace{0.5em}

\item \textbf{Correct answer: C} \\
A quota and a tariff that restrict imports by the same amount are equivalent if the \textbf{quota licenses are auctioned competitively}.  
In that case, the government captures the quota rents just as it would collect tariff revenue; domestic prices and quantities match the tariff case.

\vspace{0.5em}

\item \textbf{Correct answer: C} \\
A \textit{VER} is an export restraint negotiated by the importing country but administered by the exporting country.  
The resulting scarcity of exports raises the import price and creates rents, which are typically captured by the \textbf{exporting firms/government}, not the importer.

\vspace{0.5em}

\item \textbf{Correct answer: B} \\
A \textit{large country} can affect world prices. By restricting imports, it may reduce the world price of its imports (improving its \textbf{terms of trade}).  
This is the mechanism behind the concept of an ``optimal tariff.''

\newpage

\item \textbf{Correct answer: B} \\
The two standard deadweight loss components from a tariff are:
\begin{itemize}
    \item \textbf{Production distortion} (over-production by high-cost domestic producers),
    \item \textbf{Consumption distortion} (under-consumption by domestic consumers because of the higher price).
\end{itemize}

\vspace{0.5em}

\item \textbf{Correct answer: C} \\
The \textit{effective rate of protection} (ERP) measures the \textbf{percentage increase in domestic value added} in an industry due to the full tariff structure, including tariffs on inputs and outputs.  
It can be much higher or lower than the nominal tariff on the final good.

\end{enumerate}

\vspace{1em}
\hrule
\vspace{1em}

\section*{Problem 1 – Tariff Analysis for a Small Country (Detailed Solution)}

\textbf{Given:}
\[
Q_S = P - 50, \qquad Q_D = 400 - P, \qquad P_W = 150.
\]
$Q$ is quantity, $P$ is price in dollars. The country is \textbf{small}, so it takes the world price as given.

\subsection*{(1) Free-trade domestic price, $Q_S$, $Q_D$, and imports}

Under free trade, domestic price equals the world price:
\[
P^{FT} = P_W = 150.
\]

Domestic supply:
\[
Q_S^{FT} = P^{FT} - 50 = 150 - 50 = 100.
\]

Domestic demand:
\[
Q_D^{FT} = 400 - P^{FT} = 400 - 150 = 250.
\]

Imports under free trade:
\[
M^{FT} = Q_D^{FT} - Q_S^{FT} = 250 - 100 = 150.
\]

\[
\boxed{P^{FT} = 150,\quad Q_S^{FT} = 100,\quad Q_D^{FT} = 250,\quad M^{FT} = 150.}
\]

\subsection*{(2) Domestic price, $Q_S$, $Q_D$, and imports with a tariff $t = 30$}

A specific tariff of \$30 per unit is imposed. For a \textbf{small} country:
\[
P_W \text{ stays at } 150,\quad P^{T} = P_W + t = 150 + 30 = 180.
\]

New domestic supply at $P^{T}=180$:
\[
Q_S^{T} = P^{T} - 50 = 180 - 50 = 130.
\]

New domestic demand at $P^{T}=180$:
\[
Q_D^{T} = 400 - P^{T} = 400 - 180 = 220.
\]

Imports with tariff:
\[
M^{T} = Q_D^{T} - Q_S^{T} = 220 - 130 = 90.
\]

\[
\boxed{P^{T} = 180,\quad Q_S^{T} = 130,\quad Q_D^{T} = 220,\quad M^{T} = 90.}
\]

\subsection*{(3) Government tariff revenue}

Tariff revenue equals the tariff per unit times the number of imported units:
\[
\text{Tariff revenue} = t \cdot M^{T} = 30 \times 90 = 2{,}700.
\]

\[
\boxed{\text{Tariff revenue} = \$2{,}700.}
\]

\subsection*{(4) Deadweight loss of the tariff}

The welfare-cost triangles come from \textbf{changes} in quantities due to the higher domestic price.

\paragraph{Step 1: Quantity changes.}
\[
\Delta Q_S = Q_S^{T} - Q_S^{FT} = 130 - 100 = 30 \quad \text{(extra domestic production)}.
\]
\[
\Delta Q_D = Q_D^{FT} - Q_D^{T} = 250 - 220 = 30 \quad \text{(lost consumption)}.
\]

\paragraph{Step 2: Production distortion loss.}
On the left side of the import-demand diagram, the tariff causes overproduction by relatively high-cost domestic firms.  
The production distortion deadweight loss is the area of the left triangle:
\[
\text{DWL}_{\text{prod}} = \frac{1}{2} \times (\text{tariff}) \times (\Delta Q_S)
= \frac{1}{2} \times 30 \times 30 = 450.
\]

\paragraph{Step 3: Consumption distortion loss.}
On the right side, the higher price causes consumers to buy less, even though the world price is lower.  
The consumption distortion deadweight loss is:
\[
\text{DWL}_{\text{cons}} = \frac{1}{2} \times (\text{tariff}) \times (\Delta Q_D)
= \frac{1}{2} \times 30 \times 30 = 450.
\]

\paragraph{Step 4: Total deadweight loss.}
\[
\text{Total DWL} = \text{DWL}_{\text{prod}} + \text{DWL}_{\text{cons}} = 450 + 450 = 900.
\]

\[
\boxed{\text{Total deadweight loss} = \$900.}
\]

\paragraph{Interpretation.}
For a small country, the tariff:
\begin{itemize}
    \item Transfers some surplus from consumers to producers and government (this is a \emph{redistribution}), and
    \item Creates a net \emph{efficiency loss} of \$900 due to distorted production and consumption decisions.
\end{itemize}
There is no terms-of-trade gain to offset this loss, so national welfare is lower by \$900.

\end{document}
